\documentclass[11pt,a4paper]{article}

% ---------------------------------------------------------
% PACKAGES
% ---------------------------------------------------------
\usepackage[utf8]{inputenc}
\usepackage[T1]{fontenc}
\usepackage{lmodern}
\usepackage{amsmath, amssymb, amsfonts, bm}
\usepackage{geometry}
\usepackage{graphicx}
\usepackage{hyperref}
\usepackage{authblk}
\usepackage{cite}
\usepackage{physics}
\usepackage{mathtools}
\usepackage{tensor}
\usepackage{enumitem}

\geometry{margin=2.4cm}

% ---------------------------------------------------------
% TITLE
% ---------------------------------------------------------
\title{\textbf{Companion VII (v49.1)}\\[0.2cm]
\textbf{Boltzmann Hierarchy in a CPT-Symmetric Bipartite Universe}\\[0.2cm]
\large Foundations of Relaxation-Driven Cyclic Cosmology (RDCC)}

\author[1]{Michael Lehmann}
\affil[1]{\small Independent Researcher, Relaxation-Driven Cyclic Cosmology (RDCC)\\
\texttt{mi.lehmann@gmx.de}}

\date{September 2026}

\begin{document}
\maketitle

% ========================================================
% ABSTRACT
% ========================================================

\begin{abstract}
This Companion develops the corrected Boltzmann hierarchy of RDCC in its
Phase-1-integrated v49.1 form. The mirror sector replaces cold dark matter as a
non-compact gravitational component, eliminating the CDM perturbation equation
and introducing the mirror-sector density contrast \(\delta_{(-)}\). The Poisson
equation is formulated in fully covariant form and incorporates the
dimension-8 inter-sector operator \(T^{(+)}_{\mu\nu}T^{(-)\mu\nu}\). The
coefficients \(C_1\)--\(C_4\) are fixed theoretical factors of order unity rather
than fit parameters. The hierarchy inherits the relaxation scale
\(k_{\mathrm{rel}}=a_{\mathrm{fr}}H(a_{\mathrm{fr}})\) from Companion XII,
which separates standard and relaxation-modified perturbation growth. This
Companion provides the corrected perturbation structure that underlies RDCC's
growth suppression, tensor modulation, and survey interface.
\end{abstract}

\tableofcontents
\newpage

% ========================================================
% 1 INTRODUCTION
% ========================================================

\section{Introduction}

The Boltzmann hierarchy describes the evolution of cosmological perturbations.
In \(\Lambda\)CDM, cold dark matter (CDM) plays a central role: its density
contrast \(\delta_{\mathrm{c}}\) drives gravitational clustering and determines
the growth of structure.

RDCC replaces CDM with the mirror sector, which forms diffuse, non-compact
halos due to its reversed thermodynamic arrow. These halos contribute
gravitationally but do not undergo dissipative collapse.

In v49.1, the Boltzmann hierarchy is corrected to reflect this structure:
\begin{itemize}[leftmargin=0.8cm]
  \item the CDM perturbation equation is removed,
  \item \(\delta_{\mathrm{c}}\) is replaced by the mirror-sector contrast \(\delta_{(-)}\),
  \item the Poisson equation is fully covariant,
  \item the coefficients \(C_1\)--\(C_4\) are fixed theoretical factors,
  \item the dimension-8 operator provides the ultraviolet origin of the
        modification,
  \item the relaxation parameter \(\alpha(t)\) is globally defined and
        RG-stabilized,
  \item the freeze-out scale \(k_{\mathrm{rel}}\) of Companion XII enters the
        growth sector as a geometric transition scale.
\end{itemize}

This Companion develops the corrected Boltzmann hierarchy in a form that is
mathematically precise, physically transparent, and fully compatible with the
global structure of RDCC.

% ========================================================
% 2 COVARIANT POISSON EQUATION
% ========================================================

\section[Covariant Poisson Equation]{Covariant Poisson Equation}

The gravitational potential \(\Phi\) satisfies the covariant Poisson equation
\begin{equation}
\nabla^2 \Phi - 3H(\dot{\Phi} + H\Phi)
=
4\pi G
\left(
\rho_{(+)} \delta_{(+)}
+
\rho_{(-)} \delta_{(-)}
\right)
+
\frac{C_1}{M^4} \, \Xi,
\end{equation}
where:
\begin{itemize}[leftmargin=0.8cm]
  \item \(\delta_{(+)}\) is the visible-sector density contrast,
  \item \(\delta_{(-)}\) is the mirror-sector density contrast,
  \item \(\Xi\) encodes the geometric contribution of the dimension-8 operator,
  \item \(C_1\) is a fixed theoretical factor of order unity.
\end{itemize}

% --------------------------------------------------------
% 2.1 Elimination of CDM and Mirror-Sector Replacement
% --------------------------------------------------------

\subsection[Elimination of CDM]{Elimination of CDM}

Cold dark matter does not exist in RDCC. Its perturbation equation is removed
and replaced by the mirror-sector contrast \(\delta_{(-)}\).

The standard CDM equation
\begin{equation}
\dot{\delta}_{\mathrm{c}} = -\theta_{\mathrm{c}} + 3\dot{\Phi}
\end{equation}
is replaced by
\begin{equation}
\dot{\delta}_{(-)} = -\theta_{(-)} + 3\dot{\Phi},
\end{equation}
with \(\theta_{(-)}\) the mirror-sector velocity divergence.

The mirror sector behaves as a non-compact gravitational component.

% --------------------------------------------------------
% 2.2 Relaxation Scale and Perturbation Transition
% --------------------------------------------------------

\subsection{Relaxation Scale and Perturbation Transition}

The perturbation sector inherits a characteristic geometric scale from the
freeze-out dynamics of the relaxation sector. Companion XII established the
relaxation scale
\begin{equation}
  k_{\mathrm{rel}}
  =
  a_{\mathrm{fr}} H(a_{\mathrm{fr}}),
\end{equation}
where
\begin{equation}
  a_{\mathrm{fr}} \simeq 10^{-4}
\end{equation}
marks the freeze-out epoch of the global relaxation field.

This scale separates two perturbative regimes. For
\begin{equation}
  k < k_{\mathrm{rel}},
\end{equation}
perturbations remain largely unaffected by relaxation effects. For
\begin{equation}
  k > k_{\mathrm{rel}},
\end{equation}
relaxation-induced suppression becomes significant. The relaxation scale
therefore acts as the transition scale between standard and
relaxation-modified perturbation growth.

% ========================================================
% 3 BOLTZMANN HIERARCHY
% ========================================================

\section[Boltzmann Hierarchy]{Boltzmann Hierarchy}

The Boltzmann hierarchy for photons, baryons, neutrinos, and the mirror sector
is given by the standard equations with the corrected gravitational source
terms.

\subsection[Photon Hierarchy]{Photon Hierarchy}

\begin{align}
\dot{\delta}_{\gamma} &= -\frac{4}{3}\theta_{\gamma} + 4\dot{\Phi}, \\
\dot{\theta}_{\gamma} &= k^2 \left( \frac{\delta_{\gamma}}{4} - \sigma_{\gamma} \right)
+ k^2 \Psi.
\end{align}

\subsection[Mirror-Sector Hierarchy]{Mirror-Sector Hierarchy}

The mirror sector behaves as a non-compact gravitational component:
\begin{align}
\dot{\delta}_{(-)} &= -\theta_{(-)} + 3\dot{\Phi}, \\
\dot{\theta}_{(-)} &= -H\theta_{(-)} + k^2 \Psi.
\end{align}

% --------------------------------------------------------
% 3.1 Fixed C-Coefficients
% --------------------------------------------------------

\subsection[Fixed C coefficients]{Fixed \(C_1\)--\(C_4\) Coefficients}

The coefficients \(C_1\)--\(C_4\) are fixed theoretical factors of order unity:
\begin{align}
C_1 &\sim \mathcal{O}(1), \\
C_2 &\sim \mathcal{O}(1), \\
C_3 &\sim \mathcal{O}(1), \\
C_4 &\sim \mathcal{O}(1).
\end{align}

They are not fit parameters. Their values arise from the global geometric
structure and the dimension-8 operator.

% --------------------------------------------------------
% 3.2 Relaxation-Modified Growth
% --------------------------------------------------------

\subsection{Relaxation-Modified Growth}

The mirror-sector perturbations inherit the relaxation scale through the
freeze-out geometry. The growth factor may be written schematically as
\begin{equation}
  D_{\mathrm{RDCC}}(k,a)
  =
  D(a)
  \left[
    1
    -
    \chi_D
    \left(
      \frac{k}{k_{\mathrm{rel}}}
    \right)^{n}
  \right].
\end{equation}
Similarly, the growth rate becomes
\begin{equation}
  f_{\mathrm{RDCC}}(k,a)
  =
  f(a)
  \left[
    1
    -
    \chi_f
    \left(
      \frac{k}{k_{\mathrm{rel}}}
    \right)^{n}
  \right].
\end{equation}
The appearance of the same relaxation scale in both growth observables and
tensor damping provides evidence for a common geometric origin.

% --------------------------------------------------------
% 3.3 Connection to Background Cosmology
% --------------------------------------------------------

\subsection{Connection to Background Cosmology}

Companion XI introduced the relaxation correction to the background cosmology,
\begin{equation}
  \delta_{\mathrm{rel}}(z)
  =
  \alpha_{\mathrm{IR}}\, f_{\mathrm{rel}}(z).
\end{equation}
The perturbation hierarchy must remain consistent with this background
modification. Consequently, the growth of perturbations is influenced
indirectly through the modified expansion rate,
\begin{equation}
  H(z)
  =
  H_{\Lambda\mathrm{CDM}}(z)
  \left[
    1 + \delta_{\mathrm{rel}}(z)
  \right].
\end{equation}
The resulting coupling between background evolution and perturbation growth
provides the basis for future numerical implementations of RDCC in Boltzmann
solvers.

% --------------------------------------------------------
% 3.4 Relaxation-Dominated Regime
% --------------------------------------------------------

\subsection{Relaxation-Dominated Regime}

The perturbation spectrum naturally divides into two sectors. The large-scale
sector
\begin{equation}
  k < k_{\mathrm{rel}}
\end{equation}
remains close to the standard evolution. The small-scale sector
\begin{equation}
  k > k_{\mathrm{rel}}
\end{equation}
experiences relaxation-induced damping generated by the dimension-8 interaction
between the visible and mirror sectors. This transition represents the
perturbative counterpart to the freeze-out geometry established in
Companion XII.

% ========================================================
% 4 CONCLUSION
% ========================================================

\section{Conclusion}

This Companion develops the corrected Boltzmann hierarchy of RDCC in its
v49.1 form. The CDM perturbation equation is removed and replaced by the
mirror-sector contrast \(\delta_{(-)}\), reflecting the non-compact nature of the
mirror sector. The Poisson equation is formulated in fully covariant form and
incorporates the dimension-8 inter-sector operator.

The hierarchy inherits the geometric transition scale
\begin{equation}
  k_{\mathrm{rel}} = a_{\mathrm{fr}} H(a_{\mathrm{fr}})
\end{equation}
from Companion XII. Growth factor and growth rate are modified for
\(k>k_{\mathrm{rel}}\) and remain consistent with the background correction
of Companion XI.

The coefficients \(C_1\)--\(C_4\) are fixed theoretical factors of order unity,
ensuring that the perturbation structure is fully determined by the global
geometry. These results integrate the relaxation sector (Companion III), the
thermodynamic sector (Companion IV), the ultraviolet sector (Companion X), the
freeze-out geometry (Companion XII), and the growth suppression
(Companion XIII). They form the backbone of RDCC's perturbation theory and are
essential for the likelihood structure (Companion VIII) and parameter
estimation (Companion IX).

\end{document}